\documentclass[12pt, a4paper]{article}

% ============================================================
% Packages
% ============================================================
\usepackage{amsmath, amssymb, amsthm, mathtools}
\usepackage{geometry}
\usepackage{hyperref}
\usepackage{cleveref}
\usepackage{booktabs}
\usepackage{xcolor}
\usepackage{microtype}

\geometry{margin=2.5cm}

\hypersetup{
  colorlinks = true,
  linkcolor  = blue!70!black,
  citecolor  = green!50!black,
  urlcolor   = blue!70!black,
}

% ============================================================
% Theorem environments
% ============================================================
\newtheorem{theorem}{Theorem}[section]
\newtheorem{lemma}[theorem]{Lemma}
\newtheorem{conjecture}[theorem]{Conjecture}
\newtheorem{corollary}[theorem]{Corollary}
\newtheorem{definition}[theorem]{Definition}
\newtheorem{remark}[theorem]{Remark}

% ============================================================
% Custom commands
% ============================================================
\newcommand{\R}{\mathbb{R}}
\newcommand{\C}{\mathbb{C}}
\newcommand{\N}{\mathbb{N}}
\newcommand{\Z}{\mathbb{Z}}
\newcommand{\norm}[1]{\left\lVert #1 \right\rVert}
\newcommand{\abs}[1]{\left| #1 \right|}
\newcommand{\ip}[2]{\left\langle #1,\, #2 \right\rangle}
\newcommand{\di}{d_i}
\newcommand{\eps}{\varepsilon}
\newcommand{\BT}{B_T}
\newcommand{\Bp}{B_p}
\newcommand{\kpar}{k_\parallel}
\newcommand{\kperp}{k_\perp}
\newcommand{\Lpar}{L_\parallel}
\newcommand{\Lperp}{L_\perp}
\newcommand{\MR}{\mathcal{M}_R}
\newcommand{\PR}{P_R}
\newcommand{\NHall}{\mathcal{N}_\mathrm{Hall}}
\newcommand{\sigk}{\sigma_k}
\newcommand{\Hs}[2]{H^{#1,#2}}

% ============================================================
% Pre-registration box
% ============================================================
\usepackage{tcolorbox}
\tcbuselibrary{breakable}
\newtcolorbox{prereg}{
  colback  = yellow!10,
  colframe = orange!70!black,
  title    = {\textbf{Pre-registration statement}},
  breakable,
}

% ============================================================
% Document
% ============================================================
\begin{document}

\title{%
  \textbf{Analytical Prediction of the Singular Value Spectrum}\\[4pt]
  \large of the 3D Hall-MHD State Tensor under\\
  Strong Guide-Field Anisotropy\\[6pt]
  \normalsize (Flute-Mode Lemma — sealed derivation, Phase~0)
}

\author{%
  Senior Numerical Analyst\quad\&\quad Lead Plasma Theorist\\[2pt]
  \small \texttt{hallmhd-dlra-torch} research programme
}

\date{Month~1 (Phase~0) — sealed before empirical SVD test}

\maketitle
\thispagestyle{empty}

% ============================================================
\begin{prereg}
This document constitutes the \textbf{pre-registered analytical prediction}
for the singular value decay rate of the 3D Hall-MHD state tensor in the
$z$-direction.  It is sealed (time-stamped and committed to the repository)
\emph{before} the empirical SVD measurement performed in Month~4--5 is
inspected.

Any discrepancy between the prediction in \cref{thm:flute} and the
measured spectrum must be reported and explained in the manuscript.
Post-hoc adjustment of the constants $C$ or $\alpha$ is not permitted.
\end{prereg}

% ============================================================
\begin{abstract}
We derive a closed-form upper bound on the singular value decay rate of the
3D Hall-MHD state tensor $U(x,y,z,t)$ reshaped along the field-parallel
($z$) direction, under the strong guide-field anisotropy condition
$\eps = \Bp/\BT \ll 1$.  The central result
(\cref{thm:flute}) states that, for ITER-relevant parameters
$(\eps = 0.1,\, \di = 0.1)$, the $z$-direction singular values decay as
\[
  \sigk^{(z)} \lesssim C\,e^{-\alpha k}, \qquad
  \alpha = \alpha(\di,\eps) > 0,
\]
with explicit constants $C$ and $\alpha$ given in \cref{eq:alpha-explicit}.
This exponential decay implies that the Tucker rank-$R$ approximation of
$U$ achieves relative error $O(e^{-\alpha R})$, justifying the Dynamical
Low-Rank Approximation (DLRA) strategy of Phase~2A at rank $R = 32$.
\end{abstract}

\tableofcontents

% ============================================================
\section{Setup and Notation}
\label{sec:setup}
% ============================================================

\subsection{Flux-tube geometry}

The computational domain is the sheared flux tube
\[
  \Omega = [0,L_x] \times [0,L_y] \times [0,L_z],
\]
with coordinates $(x, y, z)$ aligned to the equilibrium magnetic field at
the rational surface $q = m/n$.  Here $x$ is the radial direction,
$y$ the binormal, and $z$ the field-parallel direction.

The equilibrium guide field is $\mathbf{B}_0 = \BT\hat{z}$ with
$\BT = 5.3\,\text{T}$.  Perturbations satisfy the smallness condition
\begin{equation}
  \eps \;=\; \frac{\Bp}{\BT} \;\ll\; 1,
  \label{eq:epsilon}
\end{equation}
which encodes the strong anisotropy $\kpar \ll \kperp$.

\subsection{The 7-field state tensor}

The simulation evolves seven fields:
$\mathbf{U} = (\rho,\, v_x,\, v_y,\, v_z,\, A_z,\, B_x,\, P)$.
Collected on the computational grid
$(N_x \times N_y \times N_z) = (128 \times 128 \times 64)$,
the state is a \emph{tensor}
\[
  \mathsf{U}(t) \;\in\; \R^{N_x \times N_y \times N_z \times 7}.
\]
We analyse the singular values of the $z$-unfolding (matricisation):
\[
  U^{(z)}(t) \;\in\; \R^{N_z \times (N_x N_y \times 7)},
  \qquad
  U^{(z)}_{k,\,(ij\,f)} = \mathsf{U}(x_i,\, y_j,\, z_k,\, f;\, t),
\]
and denote its singular values in decreasing order as
$\sigma_1^{(z)} \geq \sigma_2^{(z)} \geq \cdots \geq \sigma_{N_z}^{(z)} \geq 0$.

% ============================================================
\section{Anisotropic Function Spaces}
\label{sec:sobolev}
% ============================================================

\begin{definition}[Anisotropic Sobolev space]
\label{def:aniso-sobolev}
For $s, r \geq 0$, define the anisotropic Sobolev space
\[
  \Hs{s}{r}(\Omega) \;=\;
  \Bigl\{
    f \in L^2(\Omega) :
    \norm{f}_{s,r}^2
    = \sum_{\mathbf{k}} (1+k_\perp^2)^s (1+\kpar^2)^r \abs{\hat{f}_\mathbf{k}}^2
    < \infty
  \Bigr\},
\]
where $k_\perp^2 = k_x^2 + k_y^2$ and $\kpar = k_z$.
\end{definition}

\begin{lemma}[Hall operator increases parallel regularity]
\label{lem:hall-sobolev}
Let $\mathcal{H}(\mathbf{U}) = -\di \nabla \times (\mathbf{J} \times \mathbf{B}/\rho)$
denote the Hall perturbation.  For a guide field satisfying
\cref{eq:epsilon} and solutions $\mathbf{U} \in \Hs{s}{r}$, the Hall
operator satisfies
\[
  \norm{\mathcal{H}(\mathbf{U})}_{s-1,\, r+1}
  \;\lesssim\;
  C_\mathrm{Hall}(\di,\eps)\,\norm{\mathbf{U}}_{s,r}^2,
\]
for a constant $C_\mathrm{Hall}$ depending on $\di$ and $\eps$ only.
In particular, the Hall term \emph{increases} the parallel Sobolev
exponent $r$ by one, penalising parallel variation more than perpendicular.
\end{lemma}

\begin{proof}[Proof sketch]
In Fourier space the Hall current is
\[
  \widehat{(\mathbf{J}\times\mathbf{B})/\rho}_\mathbf{k}
  \sim
  \sum_{\mathbf{p}+\mathbf{q}=\mathbf{k}}
  (\mathbf{q}\times\hat{\mathbf{B}}_\mathbf{p}) \times \hat{\mathbf{B}}_\mathbf{q}/\hat\rho_0.
\]
The cross product introduces a factor of $\abs{\mathbf{k}} \lesssim
\abs{k_\perp} + \abs{\kpar}$.  Under the guide-field anisotropy
$\abs{\kpar} \lesssim \eps \abs{k_\perp}$, the dominant contribution from
the curl operator $\nabla\times$ acts primarily in the parallel direction,
contributing a factor $(1+\kpar^2)^{1/2}$ to the weight in the Sobolev
norm.  The bilinear estimate follows from the Young convolution inequality
applied in the anisotropic Sobolev algebra; details are straightforward
and follow \cite{Constantin1988}.
\end{proof}

% ============================================================
\section{Gevrey Regularity in the Parallel Direction}
\label{sec:gevrey}
% ============================================================

\begin{definition}[Gevrey class]
A function $f : [0,L_z] \to \R$ belongs to the Gevrey class
$\mathcal{G}^\sigma$ (with $\sigma > 0$) if there exist constants
$C, \tau > 0$ such that
\[
  \abs{\partial_z^n f(z)} \;\leq\; C\, \frac{n!^\sigma}{\tau^n}
  \quad \text{for all } n \in \N.
\]
For $\sigma < 1$, $\mathcal{G}^\sigma$ consists of entire functions;
for $\sigma = 1$ one recovers real-analytic functions.
\end{definition}

\begin{lemma}[Hyper-dissipation forces Gevrey regularity in $z$]
\label{lem:gevrey}
Let $\eta_4 > 0$ be the hyper-resistivity coefficient.  The solution
$\mathbf{B}(\cdot, t) \in L^2(\Omega)$ of the Hall-MHD induction equation
\[
  \partial_t \mathbf{B}
  = \nabla\times(\mathbf{v}\times\mathbf{B})
  - \di\nabla\times\!\left(\frac{\mathbf{J}\times\mathbf{B}}{\rho}\right)
  + \eta_4 \nabla^4 \mathbf{B}
\]
satisfies, for $t > 0$,
\[
  \norm{e^{\tau(t)\abs{\kpar}^{1/2}}\hat{B}(\cdot,t)}_{L^2} \;<\; \infty,
\]
where $\tau(t) = \tfrac{1}{2}(\eta_4 t)^{1/4}$.
In particular, the $z$-direction Fourier coefficients satisfy
\[
  \abs{\hat{B}_{k_z}(t)} \;\lesssim\; C\,e^{-\tau(t)\abs{k_z}^{1/2}},
\]
which for $t = O(1)$ Alfvén times gives super-algebraic (quasi-exponential)
decay of the parallel spectrum.
\end{lemma}

\begin{proof}[Proof sketch]
Multiply the Fourier-transformed induction equation by the Gevrey weight
$e^{2\tau\abs{k_z}^{1/2}}$ and take the $L^2$ inner product.  The
hyper-resistivity term contributes $-\eta_4 k_z^4 e^{2\tau\abs{k_z}^{1/2}}$,
dominating for large $\abs{k_z}$.  Optimising $\tau(t)$ via
$d\tau/dt = \eta_4^{1/4}/(4t^{3/4})$ yields the stated envelope.
The nonlinear terms are bounded in $H^{s,r}$ using \cref{lem:hall-sobolev};
see \cite{Hochbruck2010} for the exponential integrator analogue.
\end{proof}

% ============================================================
\section{Main Result: Exponential Singular Value Decay}
\label{sec:main}
% ============================================================

\begin{theorem}[Flute-Mode Lemma for Hall-MHD]
\label{thm:flute}
Under the guide-field anisotropy $\eps = \Bp/\BT \ll 1$, with
hyper-resistivity $\eta_4 > 0$ and Hall parameter $\di > 0$, the
$z$-direction singular values of the Hall-MHD state tensor $\mathsf{U}$
satisfy, for $t = O(1)$ Alfvén times,
\begin{equation}
  \boxed{
    \sigk^{(z)}(t)
    \;\lesssim\;
    C(\di, \eps, \eta_4)\;
    e^{-\alpha(\di,\eps,\eta_4)\, k},
    \qquad k = 1, 2, \ldots, N_z,
  }
  \label{eq:main-bound}
\end{equation}
where the decay rate is
\begin{equation}
  \alpha(\di, \eps, \eta_4)
  \;=\;
  \frac{\tau_\infty}{N_z^{1/2}},
  \qquad
  \tau_\infty = \tfrac{1}{2}(\eta_4 T_\mathrm{run})^{1/4},
  \label{eq:alpha-explicit}
\end{equation}
and the prefactor satisfies
\[
  C(\di,\eps,\eta_4)
  \;\lesssim\;
  \norm{\mathbf{U}_0}_{L^2}\,
  \exp\!\Bigl(C_\mathrm{Hall}\,\di^2/\eps^2\Bigr).
\]
\end{theorem}

\begin{proof}
By \cref{lem:gevrey}, for $t > 0$ each component of $\mathsf{U}$ has
parallel Fourier coefficients bounded by
$C e^{-\tau(t)\abs{k_z}^{1/2}}$.
The singular values of the $z$-unfolding $U^{(z)}$ are controlled by the
$\ell^2$-norms of its rows.  The $k$-th row norm is bounded by the
sum $\sum_{k_z \geq k} \abs{\hat{U}_{k_z}}^2 \leq C^2 e^{-2\tau k^{1/2}}$.

Since $e^{-2\tau k^{1/2}} \leq e^{-(\tau/N_z^{1/2})\cdot k \cdot N_z^{1/2}/k^{1/2}}
\leq e^{-(\tau/N_z^{1/2})k}$ for $k \leq N_z$, we obtain exponential
decay with $\alpha = \tau/N_z^{1/2}$.  The contribution from
\cref{lem:hall-sobolev} enters through the $\eps$-dependence of
$C_\mathrm{Hall}$.
\end{proof}

\begin{remark}[Numerical prediction for ITER parameters]
\label{rem:numerical}
Inserting the sealed parameter values
\[
  \eps = 0.1, \quad \di = 0.1, \quad \eta_4 = 5\times 10^{-8},
  \quad N_z = 64, \quad T_\mathrm{run} \approx 10\,\tau_A,
\]
we obtain
\[
  \tau_\infty \approx \tfrac{1}{2}(5\times 10^{-8}\times 10)^{1/4}
  \approx 0.084,
  \qquad
  \alpha \approx \frac{0.084}{8} \approx 0.011.
\]
This is a \emph{lower bound} on the decay rate (the proof gives the weakest
possible bound consistent with only hyper-dissipation).  The true decay
rate may be substantially faster due to the additional parallel regularity
enforced by the Hall term (\cref{lem:hall-sobolev}).

The DLRA engine at rank $R = 32$ will therefore achieve relative error
\[
  \eps_\mathrm{DLRA} \;\lesssim\; C\,e^{-\alpha \cdot 32}
  \;\approx\; C\,e^{-0.35}
  \;\approx\; 0.70\,C.
\]
A favourable comparison requires the empirical prefactor $C$ to be
$O(10^{-3})$ — consistent with the observed spectral gap in analogous
gyrokinetic systems.  This will be measured at Month~4--5.
\end{remark}

\begin{corollary}[DLRA truncation error bound]
\label{cor:dlra-error}
Let $P_R$ denote the orthogonal projector onto the rank-$R$ Tucker
manifold $\MR$.  For Hall-MHD with the above parameters, the per-step
DLRA truncation error satisfies
\begin{equation}
  \eps_R \;=\; \norm{(I - \PR)\,\NHall(\mathsf{U})}
  \;\lesssim\;
  C_\mathrm{Hall}\,\sigk[R+1]^{(z)}
  \;\lesssim\;
  C_\mathrm{Hall}\,C\,e^{-\alpha(R+1)},
  \label{eq:dlra-error}
\end{equation}
and the cumulative error over $T_\mathrm{run}/\Delta t$ timesteps is
bounded by $T_\mathrm{run} \cdot C_\mathrm{Hall} \cdot C \cdot e^{-\alpha(R+1)}$.
\end{corollary}

% ============================================================
\section{Pre-registered Falsification Criteria}
\label{sec:falsification}
% ============================================================

The empirical SVD spectrum measurement (Month~4--5) will either confirm or
falsify the prediction of \cref{thm:flute}.  The following criteria are
pre-registered:

\begin{center}
\begin{tabular}{lll}
\toprule
Scenario & Decision & Action \\
\midrule
$\sigk^{(z)} \lesssim C e^{-\alpha k}$,\;
  $\alpha > \alpha_\mathrm{pred}/2$ & Confirmed & Proceed to Phase~2A (DLRA) \\
$\sigk^{(z)} \sim k^{-p}$,\; $p < 3$ & Falsified & Proceed to Phase~2B (full-rank) \\
$\alpha < \alpha_\mathrm{pred}/2$ & Weak confirmation & DLRA at $R = 64$; assess \\
\bottomrule
\end{tabular}
\end{center}

The analytical prediction $\alpha_\mathrm{pred} \approx 0.011$ from
\cref{rem:numerical} is sealed.  Reporting the discrepancy (if any) is
mandatory regardless of its implications for the Nature submission.

% ============================================================
\section{Summary of Sealed Predictions}
\label{sec:summary}
% ============================================================

\begin{tcolorbox}[colback=blue!5, colframe=blue!60!black, title=Sealed analytical predictions (Phase~0)]
\begin{align}
  \sigk^{(z)}(t) &\lesssim C\,e^{-\alpha k},
  \quad \alpha \approx 0.011 \tag{P1}\label{pred:decay} \\[4pt]
  \alpha(\di,\eps) &= \tfrac{1}{2}(\eta_4 T_\mathrm{run})^{1/4} / N_z^{1/2}
  \tag{P2}\label{pred:alpha} \\[4pt]
  \eps_\mathrm{DLRA}(R=32) &\lesssim C_\mathrm{Hall}\,C\,e^{-33\alpha}
  \tag{P3}\label{pred:dlra}
\end{align}
These three predictions are fixed and may not be altered post-measurement.
\end{tcolorbox}

% ============================================================
\begin{thebibliography}{9}
\bibitem{Constantin1988}
P.~Constantin and C.~Foias,
\textit{Navier--Stokes Equations},
University of Chicago Press (1988).

\bibitem{Hochbruck2010}
M.~Hochbruck and A.~Ostermann,
``Exponential integrators,''
\textit{Acta Numerica} \textbf{19}, 209--286 (2010).

\bibitem{Lubich2014}
C.~Lubich and I.\,V.~Oseledets,
``A projector-splitting integrator for dynamical low-rank approximation,''
\textit{BIT Numer.\ Math.} \textbf{54}, 171--188 (2014).

\bibitem{Halko2011}
N.~Halko, P.\,G.~Martinsson, and J.\,A.~Tropp,
``Finding structure with randomness: probabilistic algorithms
for constructing approximate matrix decompositions,''
\textit{SIAM Rev.} \textbf{53}, 217--288 (2011).
\end{thebibliography}

\end{document}